\documentclass[UTF8,zihao=-4]{ctexart}

\usepackage[a4paper,top=2.2cm,bottom=2.2cm,left=2.25cm,right=2.25cm]{geometry}
\usepackage{amsmath}
\usepackage{array}
\usepackage{booktabs}
\usepackage{enumitem}
\usepackage{fancyhdr}
\usepackage{fancyvrb}
\usepackage{float}
\usepackage{hyperref}
\usepackage{longtable}
\usepackage{tabularx}
\usepackage{xcolor}

\definecolor{ScintillaBlue}{HTML}{194B70}
\definecolor{ScintillaGreen}{HTML}{176B45}
\definecolor{ScintillaRed}{HTML}{9B2C2C}
\definecolor{ScintillaGray}{HTML}{F2F4F6}

\hypersetup{
  colorlinks=true,
  linkcolor=ScintillaBlue,
  urlcolor=ScintillaBlue,
  citecolor=ScintillaBlue,
  pdftitle={ScintillaOS BuildStorm 题二 2.3 内核设计与优化实现文档},
  pdfauthor={ScintillaOS 团队}
}

\pagestyle{fancy}
\fancyhf{}
\fancyhead[L]{ScintillaOS}
\fancyhead[R]{BuildStorm 内核设计与优化实现文档}
\fancyfoot[C]{\thepage}
\setlength{\headheight}{14pt}

\setlength{\parindent}{2em}
\setlength{\parskip}{0.35em}
\setlength{\emergencystretch}{3em}
\setlength{\tabcolsep}{4pt}
\renewcommand{\arraystretch}{1.23}
\setlist{nosep,leftmargin=2.2em}
\newcolumntype{Y}{>{\raggedright\arraybackslash}X}
\newcommand{\code}{\nolinkurl}
\newcommand{\kept}{\textcolor{ScintillaGreen}{保留}}
\newcommand{\reverted}{\textcolor{ScintillaRed}{回退}}

\title{\bfseries ScintillaOS\\[0.4em]
BuildStorm 题二 2.3 内核设计与优化实现文档}
\author{雷雨泽（队长）\quad 杜宇恒\quad 宋天琦}
\date{2026 年 8 月 12 日}

\begin{document}

\maketitle
\thispagestyle{empty}

\begin{abstract}
本文说明 ScintillaOS 为全国大学生操作系统比赛 2026 年线上决赛题二 BuildStorm
所进行的兼容性开发、并发正确性修复和性能优化。目标负载不是单个定制程序，而是
在 glibc 用户态中运行真实 Rust 工具链，以多进程、多线程方式从源码构建数百个
crate，最终生成并启动 arceos-helloworld。该工作负载同时覆盖动态链接、进程和线程、
futex、epoll、虚拟内存、ext4、VirtIO 块设备、多核调度和双架构页表等内核子系统。

2026 年 8 月 12 日官网最终日志证明：RISC-V64 与 LoongArch64 均通过工具链检查、
最小项目构建、完整 release 构建、产物大小检查和产物启动检查。RISC-V64 的正式构建
耗时为 2270.92 秒，LoongArch64 为 1700.89 秒；BuildStorm 双平台原始得分合计
311.4 分，前置 CAgent 回归双平台均为 199.1 分，整次官网评测总分为 709.6 分。
本文给出从故障定位到最终结果的完整证据链，并明确区分最终保留的实现与实验后回退的
候选方案。
\end{abstract}

\vfill
\noindent\colorbox{ScintillaGray}{%
\parbox{0.96\textwidth}{
\textbf{文档对应版本。}
官网成功样本的最后关键修复对应提交 \code{266bd1aa}；其后提交
\code{8632a3e1} 只补充双架构本地构建的确定性和并行构建隔离，不冒充官网计时样本。
本文位于 \code{documents/BuildStorm题二2.3内核设计优化实现文档.tex}，
最终提交分支为 \code{finalsubmission}。
}}

\clearpage
\tableofcontents
\clearpage

\section{评审要点与结论}

\subsection{四项人工评审要求的对应位置}

\begin{table}[H]
\centering
\small
\begin{tabularx}{\textwidth}{p{0.25\textwidth}Yp{0.18\textwidth}}
\toprule
评审项 & 本文提供的证据 & 主要章节 \\
\midrule
问题定位与根因分析（6 分） & 系统调用级故障、锁等待、块设备计数、页表和 QEMU monitor 证据；区分宿主资源不足与 guest 内核错误 & 第 \ref{sec:diagnosis} 节 \\
修复或优化的设计与实现（6 分） & Linux 语义下的进程、内存、文件系统、锁、双架构和产物启动设计；完整修改覆盖矩阵 & 第 \ref{sec:design}、\ref{sec:matrix} 节 \\
实验分析（4 分） & 同条件 A/B、回退实验、完整本地长测和官网最终数据；报告耗时、请求数、锁等待和加速比 & 第 \ref{sec:experiments} 节 \\
AI 使用与可复现步骤（4 分） & AI 的真实参与边界、错误修正方式、源码和日志留痕、双架构复现命令与判定标记 & 第 \ref{sec:ai}、\ref{sec:reproduce} 节 \\
\bottomrule
\end{tabularx}
\caption{人工评审项索引}
\end{table}

\subsection{最终官网结果}

最新官网日志为 \code{Riscv输出-11.txt} 和 \code{LoongArch输出-11.txt}。
两份日志都出现 \code{TOOLCHAIN_RESULT status=OK}、
\code{MINIBUILD_RESULT status=OK} 和最终的
\code{BUILDSTORM_RESULT ... status=OK ... run=OK}，随后正常输出测试组结束标记并关机。

\begin{table}[H]
\centering
\small
\begin{tabular}{lrrrrl}
\toprule
平台 & 核数 & 正式构建耗时 & 产物大小 & 自动评测得分 & 产物启动 \\
\midrule
RISC-V64 & 8 & 2270.92 s & 1,681,000 B & 131.4 / 200 & OK \\
LoongArch64 & 12 & 1700.89 s & 1,714,568 B & 180.0 / 200 & OK \\
\midrule
双平台合计 & --- & --- & --- & 311.4 / 400 & 均通过 \\
\bottomrule
\end{tabular}
\caption{2026-08-12 官网 BuildStorm 自动评测最终结果}
\label{tab:official-result}
\end{table}

公开 judge 的 Linux 基线为 RISC-V64 的 1616.09 秒和 LoongArch64 的 1985.21 秒。
编译时间分按
\[
S_{time}=120\times\operatorname{clamp}\left(\frac{2B-t}{B},0,1\right)
\]
计算。RISC-V64 的 2270.92 秒对应 71.4 分；LoongArch64 的 1700.89 秒快于其基线，
得到 120 分。两平台还分别获得工具链 8 分、最小构建 12 分和完整构建 40 分。题 2.3
的 20 分为本文对应的人工评审项，不在上述脚本自动得分中。
同一轮 CAgent 十项在两平台均全通过且均获得 199.1 分，说明 BuildStorm 改动没有破坏
前置基础命令、网络和文件系统能力。

\section{工作负载、约束与总体架构}

\subsection{工作负载不是面向单一用例的适配}

题二首先运行 \code{rustc --version}、\code{cargo --version}，再用
\code{cargo new} 创建项目、执行真实 \code{cargo build} 并运行 Hello World；随后从零
构建 arceos-helloworld。完整构建包含 446 个左右的预构建单元和数百个 crate，持续产生
\code{fork/clone/execve/wait4}、futex、epoll、pipe、Unix socket、mmap、COW、文件创建、
rename、truncate、目录遍历和同步块 I/O。我们没有识别 crate 名称、文件名或评分脚本来
跳过工作，也没有修改 \code{/proc/uptime} 或系统时钟伪造耗时。

正式环境同时要求：保留 VirtIO 网卡；RISC-V64 和 LoongArch64 使用相同 Linux ABI；
文件系统镜像离线且可写；QEMU 为大内存多核配置；最终产物必须超过 500 KiB 并实际启动。
因此“cargo 返回 0”不是充分条件，必须保证整个产物链路可复现。

\subsection{最终数据路径}

\begin{center}
\fbox{\begin{minipage}{0.92\textwidth}
\centering
\code{bash / cargo / rustc / linker / nested QEMU}\\
$\downarrow$ Linux ABI：进程、线程、futex、epoll、socket、mmap、文件系统系统调用\\
$\downarrow$ Task / Scheduler / Signal / VM / VFS / Page Cache\\
$\downarrow$ lwext4 事务与全局哈希加 CLOCK 块缓存\\
$\downarrow$ VirtIO Block；并行保留 VirtIO Net\\
$\downarrow$ RISC-V64 Sv39 或 LoongArch64 页表、TLB 与 IPI
\end{minipage}}
\end{center}

内核把正确性边界放在 Linux ABI、地址空间和 ext4 事务层，把性能边界放在文件页缓存、
块缓存、高低水位写回、分片 allocator 和按需调度上。这样优化不是对某个测试命令调参，
而是降低真实复杂软件构建的通用成本。

\section{问题定位与根因分析}
\label{sec:diagnosis}

\subsection{定位方法：先短测，再长测，最后官网闭环}

早期单次完整构建需要数小时，直接长测会把定位周期拉到不可接受。我们建立了四层证据：
第一层是 \code{cargo check}、双架构内核构建和 ABI 静态审计；第二层是 syscall 语义小测、
\code{FS_QUICK}、COW、VM、SMP、公平性和嵌套 QEMU 启动测试；第三层是真实官方磁盘的
工具链、minibuild、446 单元预构建和 release 构建；第四层才是官网 judge 输出。

诊断构建记录页表锁、帧分配器、Rust heap、调度器、VFS 和块设备的调用、竞争、累计等待、
最长等待、DMA 字节数与请求大小。长时间停顿时通过 QEMU monitor 的 CPU 快照判断所有 hart
是在用户态计算、设备等待、锁等待还是 guest panic。宿主 Docker/Rosetta 的资源终止不被
当作内核语义失败；只有新鲜覆盖层上可重复的 guest 证据才进入修复结论。

\subsection{启动速度和时间源的根因}

最初内核把实时时钟偏移混入运行时间，\code{/proc/uptime} 也缺少稳定实现。BuildStorm
以 guest uptime 计时，工具链内部的超时又依赖单调时钟；两者混用会使墙钟校准影响超时和
评分。另一个大头是启动阶段直接触碰或清零大块物理内存，并在 \code{execve} 时立即读取
全部 ELF 段和分配完整用户栈。Rust 工具链频繁启动短命令，这些工作大部分从未被访问。

\subsection{多线程进程语义缺口}

真实 glibc 和 Rust runtime 会组合使用 \code{CLONE_VM}、\code{CLONE_FILES}、
\code{CLONE_THREAD}、\code{clear_child_tid}、futex requeue、\code{FD_CLOEXEC} 和 wait。
旧实现把“线程 ID”“进程 ID”“子进程所有者”和“地址空间所有者”混在一起，导致：
同一线程组中的非 leader 无法等待组内创建的子进程；\code{execve} 后共享 fd 表生命周期
错误；退出线程没有正确清零并唤醒 \code{clear_child_tid}；信号唤醒与 Ready--Running
交接之间还可能把正在运行的任务重复入队。

futex 的旧 key 只使用用户虚拟地址。不同进程恰好使用相同虚拟地址时会错误互相唤醒；
反过来，同一 \code{MAP_SHARED} 页在不同地址映射时又不能共享等待队列。requeue 或超时后，
任务还可能残留在旧 key 中，最终使 cargo 永久等待。首次完整 RV 构建在 \code{0/446}
出现 \code{pid ... was enqueued twice}，正是信号提升优先级时未重新核对任务状态的竞态。

\subsection{用户内存访问与 COW 竞态}

旧的 syscall 写回路径先检查用户页可写，再释放 VM 锁执行实际复制。另一个线程可在两步
之间 \code{fork}，把同一 PTE 改成只读 COW；内核随后写入旧翻译结果会得到 \code{EFAULT}，
表现为 cargo 指纹文件、\code{execve} 参数或 socket 地址随机出现 \code{Bad address}。
这不是简单增加一次重试可以完全解决的问题，因为校验、触发 COW 和复制必须属于同一个
地址空间同步协议。

\code{MADV_DONTNEED} 原先只改变 VMA 元数据，没有移除 PTE 和释放物理页；8 GiB 环境下
rustc 反复映射和丢弃内存，最终形成确定性的常驻页增长。共享 VM 的多核任务还要求在
\code{mprotect}、\code{munmap}、COW 和 \code{mremap} 后向所有正在运行该地址空间的 CPU
发送 TLB shootdown，否则远端 hart 可继续使用旧写权限或已释放物理页。

\subsection{文件系统是主要性能瓶颈}

初始 ext4 路径的三个大头是：每个 4 KiB miss 都线性扫描固定块缓存；脏块逐块同步写回；
每次 \code{stat/open/lookup} 都重复走完整路径并进入 lwext4 全卷锁。Rust 构建包含大量
小文件、元数据查询、临时文件 rename 和顺序读取，导致 CPU 数增加后反而在同一把卷锁和
VirtIO 队列前排队。

语义上也存在会中断构建的问题。目录 rename 的链接计数按“新增目录”处理会破坏父目录
计数；\code{ftruncate} 后的 \code{stat} 错把当前文件偏移当作文件大小；同一 inode 的
多个打开实例没有共享最新 size；稀疏写、打开后 unlink、rename 替换和目录收缩后的迭代器
EOF 边界不完整。只优化缓存而不修这些语义，会更快地产生错误产物。

\subsection{SMP 锁与内存分配瓶颈}

内核早期从单核 \code{UPSafeCell} 迁移到 SMP 锁后，仍存在两个结构性问题。一是所有
动态元数据都由单一 Rust buddy heap 管理，四路 rustc 下该锁累计等待超过 13 秒；大量
进程同时退出时集中释放页表和 VMA，1 GiB 压力段需要 30.98 秒。二是睡眠读写锁的等待条件
复查调用 \code{try_read/try_write}，它会临时获取后释放锁并改变唤醒序列，却没有对应通知，
LoongArch 预构建稳定停在 \code{438--445/446}。

仅把内存序从 Acquire/Release 加强到 SeqCst 没有解决问题：故障点只从 441 移到 443，
随后又在 444 重现。QEMU monitor 显示底层锁状态已经空闲，但仍有 4 个等待者，因此根因
不是“写入不可见”，而是用于判断条件的函数本身有副作用并造成丢唤醒。

\subsection{LoongArch 架构特有根因}

LoongArch 最初只在 RISC-V 路径上完整验证。多核重负载暴露出页表 MAT、NR/NX 属性编码
不完整，以及 TLB refill 对空 PTE 的处理错误；这会把普通按需缺页变成无法恢复的 refill
循环。LoongArch 的 1 ms 定时器若每次都轮询 smoltcp 和扫描所有 interval timer，也会
产生高频全局锁竞争。其 VirtIO-Net 又使用 PCI 设备而非 RISC-V 的 MMIO 设备，需要正确
枚举 BAR 并复用同一网络状态机。

\subsection{编译成功后仍可能丢失全部分数}

首次官网样本曾完成编译，却没有通过最终产物运行。嵌套 QEMU 和 shell 会对 stdin/stdout、
Unix socket 与网络 socket 执行 \code{F_SETFL(O_NONBLOCK)}、\code{sendmsg/recvmsg}、
\code{sendto/recvfrom}。旧 File trait 对不支持状态修改的描述符返回失败，使子 QEMU 启动
前退出；socket 路径又直接解引用用户指针，在 COW 或跨页情况下可能返回错误。

LoongArch glibc 还根据 \code{AT_HWCAP} 判断非对齐访问能力。没有报告
\code{HWCAP_LOONGARCH_UAL} 时，最终静态产物在启动阶段走入不兼容路径。因此必须把
“编译返回 0、产物存在、产物大小、目标架构和产物实际启动”作为一条不可拆分的验收链。

\section{修复与优化设计}
\label{sec:design}

\subsection{时间、计时与热中断路径}

内核新增独立单调时间源，\code{CLOCK_MONOTONIC} 和 \code{/proc/uptime} 只使用自启动以来
的单调计数；\code{CLOCK_REALTIME} 使用单独的 epoch offset。实时时钟校准不再改变 sleep、
futex timeout 或 BuildStorm 计时。

LoongArch 保留 1 ms 定时精度，但网络轮询只在 10 ms 调度 tick 或设备需要时进行；smoltcp
时间戳来自同一单调源。interval timer 维护全局最近截止时间的原子 hint，只有当前时间达到
该 hint 才扫描任务并发布新的最早截止时间，避免每次中断遍历全部进程。

\subsection{进程、线程、fd 与等待语义}

任务结构显式维护稳定的 TGID、共享地址空间对象和共享 \code{files_struct}。实现
\code{CLONE_FILES} 后，fd 表的引用计数、\code{FD_CLOEXEC}、\code{close_range} 和
\code{execve} 关闭协议以共享对象为边界；\code{execve} 前先收敛同线程组其他线程，再替换
地址空间。pipe 改为 64 KiB 环形缓冲，并先在内核 scratch 中完成用户复制，避免持设备锁
触发缺页。

\code{wait4/waitid} 按线程组扫描和回收子进程，子进程退出时唤醒整个父线程组中可能等待的
线程。signal 发送遵守 \code{SchedulerLock -> TaskLock} 锁序，在提升 Ready 任务优先级前
重新核对状态；若目标已在另一个 CPU 上 Running，只发送重调度 IPI，不再建立第二个队列
所有权。重复入队断言被保留，用于捕获新的调度器错误。

\subsection{futex、epoll 与 IPC}

私有 futex key 使用 \code{address_space_id + page/offset}，只在明确的
\code{MAP_SHARED} 映射上使用共享后备对象标识。wait、wake、requeue、超时和信号中断都从
所有旧 key 清理节点；线程退出按 Linux 语义写零 \code{clear_child_tid} 并 futex wake。

RISC-V64 与 LoongArch64 的 \code{epoll_event} 按 Linux UAPI 使用 16 字节布局，events
位于偏移 0，data 位于偏移 8。EPOLLET 不再只比较两次瞬时电平，而为 file/socket 维护
\code{readiness_generation}，事件在两次采样之间出现又消失时仍不会丢边沿。内核同时补齐
\code{eventfd}、AF_UNIX \code{socketpair}、\code{sendmsg/recvmsg} 以及必要的 poll 行为，
为 cargo jobserver、rustc 和嵌套 QEMU 提供通用 IPC。

\subsection{虚拟内存、按需装入与 COW}

用户访问改为“在 VM 协议下准备页面，再复制，发现映射代际改变则重新准备”的统一路径。
内核写用户页时可在持有地址空间写同步的情况下触发 COW；\code{fork}、\code{mprotect} 和
\code{munmap} 不能插入校验与复制之间。socket、pipe、wait、stat 和 exec 参数都复用该
安全复制接口，而不是直接解引用用户指针。

\code{MADV_DONTNEED} 依次移除 PTE、收集活动 CPU、执行远端 TLB shootdown，再释放物理帧；
\code{munmap/mprotect/mremap} 使用相同活动地址空间协议。resident frame map 分片，减少
多核缺页和解除映射时对单一 BTreeMap 的争用。

ELF \code{PT_LOAD} 段和匿名用户栈改为按需分配。第一次缺页按 VMA 权限、文件偏移、非页
对齐头尾和 BSS 零填充规则加载；只读/可执行文件页进入全局文件页缓存，私有可写页保持
进程私有并沿现有 COW 处理。启动阶段也不再提前清零全部 8 GiB 物理内存，只在帧真正分配
给用户前完成零化。该设计与 Linux demand paging 一致，不减少程序可见数据。

\subsection{VFS 语义与缓存一致性}

同一文件的 size、代际和脏页状态以 superblock 与 inode 号为 key 共享，解决多打开实例、
稀疏写和截断后的可见性。inode 元数据通过单次 raw snapshot 取得，避免 \code{stat} 为
mode、uid、gid、size 和时间戳重复遍历路径。正 dentry、负 dentry、inode snapshot 和
可复用句柄模板都有代际；create、unlink、rename、truncate 或 writeback 后按父目录 inode
和目标 inode 精确失效，不能把陈旧负缓存当作 ENOENT。

create/remove/rename 以父目录 inode 为原子命名空间锁域，并在 lwext4 事务内再次核对源、
目标和 expected inode，消除 TOCTOU。rename 目录时不再按新建目录重复增加链接计数；
\code{ftruncate} 更新 inode size 而非文件 offset；打开后 unlink 的对象由 inode 引用继续
存活；目录缩短后迭代器把空项正确视作 EOF，避免在 \code{readdir} 中断言或越界。

\subsection{文件页缓存、块缓存与 I/O 合并}

文件读取增加按 inode 分片的页缓存和固定 64 KiB 顺序预读。相同页的并发 miss 通过
in-flight 状态合并，只由一个线程读取，其他线程等待代际完成；truncate、写入和回收会使
旧代际失效。写路径使用按 inode 共享的延迟写回页，读可以直接命中尚未落盘的最新脏数据。

块设备层把线性槽位查找改为预分配哈希索引，并使用全局 CLOCK 淘汰近似 LRU。脏块达到
高水位时批量写到低水位，连续物理块最多合并为 128 KiB 的逻辑写回，再按驱动 64 KiB
bounce 边界提交。高低水位滞回避免每个 miss 都触发一轮写回；写回前后都检查块代际，
保持 journal 和并发读写一致性。

\subsection{allocator、调度器与锁}

物理帧 allocator 和动态内核 heap 按平台 CPU 数分片。分配优先当前 CPU 的 shard，耗尽时
再尝试其他 shard；释放根据指针或物理地址回到原 shard，而不是按当前 CPU 回收，因此任务
迁移和跨核析构不会破坏 buddy 边界。8 B 到 4 KiB 的小对象增加每 CPU 规格化空闲链表，
缓存溢出仍按地址归还原始 shard。固定 64 MiB 启动堆作为早期启动和兜底；大内存机器从
物理内存追加按 CPU 切分的动态区，增长上限为 960 MiB。在 8 GiB guest 中该上限仍保证
至少约 7 GiB 可用于用户页和文件页缓存。

周期性时间片到期时参考 Linux \code{need_resched} 思路：若本 CPU 没有本地竞争任务、亲和性
仍合法且调度器状态可无阻塞读取，就继续运行当前任务；存在竞争、主动 yield、阻塞、退出、
亲和性变化或锁读取失败时仍调度。睡眠读写锁的等待谓词只读取
\code{reader_count/writer_count}，不再调用有副作用的 \code{try_read/try_write}；状态发布、
等待者登记和唤醒顺序保持明确的 acquire/release 关系。

\subsection{LoongArch 页表、TLB 与设备}

LoongArch PTE 按权限设置 MAT、NR 和 NX 属性，TLB refill 遇到无效或空 PTE 时进入统一缺页
路径，而不是装入错误映射或自旋。共享地址空间维护 active CPU mask，修改映射后用序列号、
IPI 和 ACK 完成远端全 TLB shootdown，空闲 hart 也能响应。PCI VirtIO-Net 完成设备枚举、
BAR 映射和中断/轮询接入，继续复用 RISC-V 的 smoltcp socket 状态机。

\subsection{最终产物启动与确定性构建}

默认 File trait 的 \code{set_flags} 对无法改变状态但允许忽略的描述符按 Linux 兼容行为返回
成功，真正支持非阻塞的 terminal、pipe、eventfd、Unix socket 和网络 socket 保存并应用
\code{O_NONBLOCK}。\code{sendto/recvfrom/sendmsg/recvmsg} 使用安全用户复制并保持消息边界。
LoongArch auxv 报告 \code{HWCAP_LOONGARCH_UAL}，使 glibc 与实际硬件非对齐访问能力一致。

构建系统为 RISC-V64 和 LoongArch64 分离 lwext4 C 产物和用户态输出，顶层双架构
\code{make all} 使用确定顺序，避免并发 Make 复用错误架构的静态库。该增量位于
\code{8632a3e1}，不改变官网成功样本的内核运行语义，只降低最终打包时的交叉架构污染风险。

\section{实验分析}
\label{sec:experiments}

\subsection{实验口径}

所有 A/B 都使用相同架构、核数、内存、磁盘镜像和工作量，相邻运行以减少宿主波动。
负数代表耗时、请求或等待下降。Mac 本地的 x86\_64 QEMU 经 Rosetta 运行，绝对 wall-clock
不能外推到官网原生评测机；因此本地表主要用来证明方向和排除回归，最终时间结论只采用
官网 guest \code{/proc/uptime}。

\subsection{最终保留优化的逐项 A/B}

\begin{longtable}{p{0.27\textwidth}p{0.20\textwidth}p{0.14\textwidth}p{0.14\textwidth}p{0.14\textwidth}}
\caption{最终保留优化的独立实验}\label{tab:ab-kept}\\
\toprule
修改 & 负载/指标 & 修改前 & 修改后 & 变化 \\
\midrule
\endfirsthead
\toprule
修改 & 负载/指标 & 修改前 & 修改后 & 变化 \\
\midrule
\endhead
全局哈希 CLOCK 与合并写回 & \code{FS_QUICK} 耗时 & 71.61 s & 44.30 s & $-38.14\%$ \\
同上 & 块写请求 & 16,926 & 2,475 & $-85.38\%$ \\
同上 & 写 DMA 时间 & 12.818809 s & 1.757420 s & $-86.29\%$ \\
固定 64 KiB 顺序预读 & \code{FS_QUICK} 耗时 & 56.47 s & 39.07 s & $-30.81\%$ \\
同上 & 读请求数 & 15,115 & 3,465 & $-77.08\%$ \\
inode 共享脏读缓存 & 512 文件完整读取 & 5,130 ms & 946 ms & $-81.56\%$ \\
VFS 正负缓存和句柄模板 & \code{FS_QUICK} 耗时 & 39.17 s & 34.70 s & $-11.41\%$ \\
同上 & raw inode 查询 & 5,915 & 383 & $-93.52\%$ \\
冷目录缓存与代际 & 2,500 文件扫描 & 15,362 ms & 6,095 ms & $-60.32\%$ \\
匿名栈按需分配 & 四路 rustc 并发段 & 389.36 s & 377.78 s & $-2.97\%$ \\
ELF \code{PT_LOAD} 按需装入 & \code{cargo --version} 中位数 & 9.07 s & 6.36 s & $-29.88\%$ \\
同上 & 完整 2.1 cargo build & 111.24 s & 102.16 s & $-8.16\%$ \\
每 CPU 小对象缓存 & 四路 rustc 并发段 & 126.95 s & 125.54 s & $-1.11\%$ \\
同上 & buddy 锁获取 & 168,379 & 33,448 & $-80.14\%$ \\
动态 heap 分片 & \code{FS_QUICK} 耗时 & 35.45 s & 32.39 s & $-8.63\%$ \\
同上 & 1 GiB 并行退出 & 30.98 s & 8.18 s & $-73.60\%$ \\
同上 & buddy 累计等待 & 145.893205 s & 0.161926 s & $-99.89\%$ \\
按需周期调度 & 调度器锁调用 & 235,213 & 169,668 & $-27.87\%$ \\
\bottomrule
\end{longtable}

块缓存优化后，平均写请求从约 8,977 字节增至 62,052 字节，约为 6.91 倍；脏块淘汰次数
从 14,095 降为 0。inode-directed rename 使局部持锁时间下降 42.73\%，卷锁等待下降
18.62\%，但整轮 wall-clock 基本不变，因此只把它作为锁粒度收益，不夸大整体加速。

\subsection{并发故障的阶段性证据}

\begin{table}[H]
\centering
\small
\begin{tabularx}{\textwidth}{p{0.24\textwidth}p{0.21\textwidth}Yp{0.16\textwidth}}
\toprule
故障/候选 & 修改前表现 & 修复后的证据 & 结论 \\
\midrule
信号优先唤醒竞态 & RV 在 \code{0/446} 重复入队 panic & 重新核对 Ready/Running 后越过原点至 \code{41/446}，后续完整通过 & 确定修复 \\
256 MiB heap 耗尽 & \code{41/446} 出现 lwext4 malloc failed 与 OOM & 扩容、分片与释放路径修复后完成全部 446 单元 & 确定修复 \\
VM 锁 Acquire/Release & LA 从 \code{438--441/446} 移到 \code{443/446} 仍停滞 & 仅移动故障窗口 & \reverted \\
VM 锁全部 SeqCst & LA 越过 443 后在 \code{444/446} 重现空锁等待者 & 排除单纯内存序假设 & \reverted \\
无副作用等待谓词 & 空闲锁仍有等待者 & 使用 reader/writer count 后自然越过 \code{438--445/446}，预构建与正式构建完成 & \kept \\
目录迭代器收缩 & LoongArch readdir panic & 空目录项按 EOF 处理后 446 单元预构建完成 & \kept \\
\bottomrule
\end{tabularx}
\caption{从可重复故障到完整构建的证据链}
\end{table}

\subsection{反向实验与回退决策}

性能工作中不以“锁次数更少”代替真实结果。以下候选均通过基本语义检查，但没有进入最终
源码；记录它们是为了说明设计选择经过反证，而不是只保留有利样本。

\begin{longtable}{p{0.29\textwidth}p{0.25\textwidth}p{0.25\textwidth}p{0.12\textwidth}}
\caption{未保留候选及回退原因}\label{tab:reverted}\\
\toprule
候选 & A/B 结果 & 根因判断 & 状态 \\
\midrule
\endfirsthead
\toprule
候选 & A/B 结果 & 根因判断 & 状态 \\
\midrule
\endhead
VirtIO bounce 64 KiB 扩至 128 KiB & \code{FS_QUICK} 34.77 s 增至 38.03 s（$+9.38\%$） & 大缓冲复制成本增加约 175\%，抵消请求合并 & \reverted \\
顺序预读自适应扩至 128 KiB & 相同工作量无稳定收益 & 短 ELF 和离散元数据产生无效预读 & \reverted \\
文件缺页 fault-around 聚簇映射 & cargo 6.20 s，处于波动；页表锁调用 $+63.21\%$ & 预读命中不代表进程会访问相邻页 & \reverted \\
私有可写 ELF 页初始共享 & cargo $+5.17\%$；四路 rustc $+1.47\%$ & 很快发生 COW，新增写缺页和页表修改 & \reverted \\
目录删除查询结果复用 & 两轮均值 35.85 s 增至 38.55 s（$+7.52\%$） & 局部查询占比小，事务内复核仍不可省 & \reverted \\
allocator 8/32 批处理 & 125.54 s 增至 126.35 s（$+0.65\%$） & buddy 合并延长关中断临界区 & \reverted \\
allocator 2/8 批处理 & 125.54 s 增至 127.19 s（$+1.31\%$） & 锁次数下降但累计等待上升 & \reverted \\
文件缓存 ghost admission & 双 rustc 83.09 s 增至 88.39 s & 一次性数据多于短期重复访问 & \reverted \\
\bottomrule
\end{longtable}

\subsection{完整长测与官网结果}

本地干净覆盖层长测中，RISC-V64 的 446 单元预构建从 137 分 47 秒降至 120 分 24 秒，
约缩短 12.6\%；其正式 release 构建耗时 3459.62 秒，产物 1,681,000 字节。LoongArch64
预构建为 54 分 51 秒，正式构建 1642.49 秒，产物 1,714,568 字节。由于本地宿主并发负载
和 Rosetta 影响，这些数字用于证明全链路完成，不作为官网基线。

官网最终结果如表 \ref{tab:official-result}。RISC-V64 日志在 36 分 59 秒完成 release profile，
最终 guest 计时为 2270.92 秒；LoongArch64 在 27 分 40 秒完成，最终计时 1700.89 秒。
两边都打印对应目标三元组的产物路径、非零字节数和 \code{run=OK}，随后测试组结束。

\begin{Verbatim}[fontsize=\small,frame=single]
BUILDSTORM_RESULT mode=multi status=OK rc=0 cores=8 \
  elapsed_s=2270.92 artifact=target/riscv64gc-unknown-linux-musl/
  release/arceos-helloworld bytes=1681000 run=OK

BUILDSTORM_RESULT mode=multi status=OK rc=0 cores=12 \
  elapsed_s=1700.89 artifact=target/loongarch64-unknown-linux-musl/
  release/arceos-helloworld bytes=1714568 run=OK
\end{Verbatim}

\section{题 2.2 修改覆盖矩阵}
\label{sec:matrix}

本节把开发记录、git 历史和最终源码逐项对齐。表中的“覆盖”表示该能力已进入最终
\code{finalsubmission}；“实验记录”表示只用于诊断且已回退，不计入最终实现。

\begin{longtable}{p{0.19\textwidth}p{0.31\textwidth}p{0.29\textwidth}p{0.11\textwidth}}
\caption{BuildStorm 2.2 全部修改覆盖审计}\label{tab:coverage}\\
\toprule
子系统 & 最终实现与 Linux 语义 & 主要源码位置 & 状态 \\
\midrule
\endfirsthead
\toprule
子系统 & 最终实现与 Linux 语义 & 主要源码位置 & 状态 \\
\midrule
\endhead
评测入口 & CAgent 后运行完整 2.1+2.2；双盘 chroot；保留网卡；结束后关机 & \code{user/src/bin/initproc.rs}, \code{Makefile} & 覆盖 \\
计时 & monotonic/realtime 分离；真实 \code{/proc/uptime}；不修改评分时间 & \code{os/src/time.rs}, \code{os/src/fs/proc.rs} & 覆盖 \\
动态装载 & glibc \code{PT_INTERP}、PIE/TLS、auxv、相对路径、symlink 与解释器路径 & \code{os/src/mm/memory_set.rs}, \code{os/src/task/task.rs}, \code{os/src/syscall/process.rs} & 覆盖 \\
按需 exec & ELF \code{PT_LOAD} 与用户栈 demand paging；BSS 和非页对齐边界 & \code{os/src/mm/memory_set.rs}, \code{os/src/task/task.rs} & 覆盖 \\
clone/exec & CLONE\_VM、CLONE\_FILES、线程组与 exec 线程收敛 & \code{os/src/task/task.rs}, \code{os/src/syscall/process.rs} & 覆盖 \\
fd 生命周期 & 共享 files\_struct、FD\_CLOEXEC、close\_range、dup/fcntl & \code{os/src/task/task.rs}, \code{os/src/syscall/fs.rs} & 覆盖 \\
pipe & 64 KiB 环形缓冲、阻塞/非阻塞、锁外用户复制 & \code{os/src/fs/pipe.rs}, \code{os/src/syscall/fs.rs} & 覆盖 \\
futex & 私有地址空间 key、MAP\_SHARED key、requeue、timeout 清理、clear\_child\_tid & \code{os/src/syscall/process.rs}, \code{os/src/task} & 覆盖 \\
等待与退出 & wait4/waitid 线程组子进程范围；父线程组唤醒；退出资源顺序 & \code{os/src/task/mod.rs}, \code{os/src/syscall/process.rs} & 覆盖 \\
signal/调度 & Ready--Running 原子复查、IPI 重调度、无竞争时跳过无效周期切换 & \code{os/src/task/manager.rs}, \code{os/src/task/processor.rs} & 覆盖 \\
epoll/eventfd & 16 字节 epoll\_event ABI、EPOLLET generation、eventfd 与 poll & \code{os/src/fs/epoll.rs}, \code{os/src/fs/eventfd.rs} & 覆盖 \\
Unix socket & socketpair、消息边界、sendmsg/recvmsg、非阻塞语义 & \code{os/src/fs/unix_socket.rs}, \code{os/src/net/syscall.rs} & 覆盖 \\
网络 socket & sendto/recvfrom 安全用户复制；双架构 VirtIO-Net；smoltcp 单调时钟 & \code{os/src/net}, \code{os/src/drivers/net} & 覆盖 \\
用户内存 & VM 同步下 prepare/copy/retry，消除 fork COW 与 syscall 写回竞态 & \code{os/src/mm/memory_set.rs}, \code{os/src/syscall} & 覆盖 \\
COW 与共享 VM & fork COW、活动 CPU mask、远端 TLB shootdown 与 ACK & \code{os/src/mm/memory_set.rs}, \code{os/src/task/processor.rs} & 覆盖 \\
mmap 系列 & mmap/munmap/mprotect/mremap；MADV\_DONTNEED 真正撤 PTE 和回收页 & \code{os/src/syscall/process.rs}, \code{os/src/mm/memory_set.rs} & 覆盖 \\
文件页缓存 & 只读页共享、in-flight 合并、64 KiB 预读、代际失效 & \code{os/src/mm/file_page_cache.rs} & 覆盖 \\
ext4 基本语义 & 稀疏文件、共享 inode size、打开后 unlink、truncate/stat 正确大小 & \code{os/src/fs/inode.rs}, \code{os/src/fs/vfs.rs}, \code{lwext4_rust} & 覆盖 \\
目录命名空间 & create/remove/rename 原子化；目录链接计数；替换和 EOF 边界 & \code{os/src/fs/vfs.rs}, \code{lwext4_rust/c/lwext4/src/ext4.c} & 覆盖 \\
VFS 缓存 & 正/负 dentry、inode snapshot、句柄模板、父目录和 inode 代际 & \code{os/src/fs/vfs.rs}, \code{os/src/fs/inode.rs} & 覆盖 \\
延迟写回 & 按 inode 共享脏页、最新脏读、高低水位 flush 与一致性失效 & \code{os/src/fs/inode.rs}, \code{os/src/fs/block_device_stream.rs} & 覆盖 \\
块缓存 & 哈希索引、全局 CLOCK、连续块合并、64 KiB 设备边界 & \code{os/src/fs/block_device_stream.rs} & 覆盖 \\
帧 allocator & 地址分片、按原 shard 回收、按需清零 & \code{os/src/mm/frame_allocator.rs} & 覆盖 \\
内核 heap & 每 CPU 小对象缓存、64 MiB 启动堆、最高 960 MiB 动态 shard & \code{os/src/mm/heap_allocator.rs}, \code{os/src/config.rs} & 覆盖 \\
VM 锁 & 无副作用 reader/writer count 等待谓词，修复 LoongArch 丢唤醒 & \code{os/src/sync/lock.rs}, \code{vendor/polyhal/src/utils/mutex_no_irq.rs} & 覆盖 \\
热中断 & 网络 10 ms 节拍、interval timer 最近 deadline hint & \code{os/src/main.rs}, \code{os/src/task/manager.rs} & 覆盖 \\
LoongArch MMU & PTE MAT/NR/NX、空项 TLB refill、IPI TLB 刷新 & \code{vendor/polyhal/src/pagetable/loongarch64.rs}, \code{vendor/polyhal-trap} & 覆盖 \\
LoongArch ABI & \code{AT_HWCAP=HWCAP_LOONGARCH_UAL} 与 glibc 启动一致 & \code{os/src/task/task.rs} & 覆盖 \\
终端与产物 & F\_SETFL/O\_NONBLOCK、默认 File set\_flags、嵌套 QEMU socket 用户复制 & \code{os/src/fs/stdio.rs}, \code{os/src/fs/mod.rs}, \code{os/src/net/syscall.rs} & 覆盖 \\
双架构构建 & lwext4 和 user 产物按 target 隔离，顶层确定顺序打包 & \code{lwext4_rust/build.rs}, \code{lwext4_rust/c/lwext4/Makefile}, \code{user/Makefile}, \code{Makefile} & 覆盖 \\
反向候选 & 大 bounce、自适应预读、fault-around、可写 ELF 共享、批量 allocator 等 & 表 \ref{tab:reverted} 与 git/开发记录 & 实验记录 \\
\bottomrule
\end{longtable}

覆盖矩阵包括从工具链启动到最终产物运行所需的全部 2.2 增量，也包括性能实验中确认无效并
回退的分支。它不把更早的初赛 LTP、lmbench 或面向其他测试的修改重复包装为 BuildStorm
优化；只有直接支撑复杂构建或保证其前置回归的内容进入本文。

\subsection{关键提交与证据映射}

\begin{longtable}{p{0.16\textwidth}p{0.31\textwidth}p{0.42\textwidth}}
\caption{关键提交的功能边界}\label{tab:commits}\\
\toprule
提交 & 主要阶段 & 与题 2.2 的关系 \\
\midrule
\endfirsthead
\toprule
提交 & 主要阶段 & 与题 2.2 的关系 \\
\midrule
\endhead
\code{df511f76} & 八核 BuildStorm/Q2.1 基础 & epoll、eventfd、Unix socket、共享 fd、SMP 调度、COW、TLB 和阻塞唤醒 \\
\code{ecc1a5fe} & 双架构迁移 & 合并主线并补齐 LoongArch VirtIO-Net 与双架构评测入口 \\
\code{6ec1e609}、\code{b469dbd1} & 评分协议与并发回归 & 测试组标记、网络并发、调度上下文和阻塞唤醒 \\
\code{90803c5e} & 完整构建语义 & 时间源、futex、用户复制、epoll ABI、wait 线程组及 ext4 语义 \\
\code{01c7091b} & 文件系统大头优化 & VFS 缓存、文件页缓存、稀疏/共享 inode、命名空间锁和 I/O 热路径 \\
\code{dae87457} & 内存与完整长测稳定性 & CLOCK 块缓存、合并写回、按需 ELF/栈、allocator 分片、调度竞态和大内存容量 \\
\code{d293494c} & LoongArch MMU & PTE 属性与 TLB refill 卡死修复 \\
\code{881565f2} & 产物终端 & terminal 和描述符非阻塞语义 \\
\code{09642b80} & LoongArch 长测收口 & VM 锁丢唤醒、目录迭代 EOF 和网络轮询修复 \\
\code{18d0c340} & 嵌套 QEMU socket & send/recv 用户内存安全复制与消息边界 \\
\code{266bd1aa} & 官网成功版本的最后关键修复 & 默认 File flags 与 LoongArch UAL HWCAP，双架构产物均可启动 \\
\code{8632a3e1} & 发布构建确定性 & 按 target 隔离 lwext4/user 产物并固定双架构构建顺序；晚于官网成功计时 \\
\bottomrule
\end{longtable}

\section{AI 辅助使用说明}
\label{sec:ai}

本项目使用 OpenAI Codex 系列大模型和 AI Agent 作为开发辅助工具。AI 的实际参与并不限于
文字润色：它协助检索 Linux/UAPI 语义，阅读仓库和日志，生成过部分 Rust/C 修改草案、
诊断插桩、压力测试和文档初稿，并在长测期间帮助归纳锁、页表、系统调用与文件系统之间的
关联。与 AI 的交流记录、每次命令输出、git 提交、开发文档和官网日志均被保留，可用于
追溯哪些建议进入过候选代码。

核心目标、总体架构取舍和验收标准由队员掌握并决定。团队明确要求修复必须符合 Linux
通用语义，不能按 crate 名称或评测输出绕过工作；队员决定测试优先级、分支发布时机、
双架构和回归边界，审查 AI 给出的补丁，并通过本地 QEMU、Docker、官方镜像和官网评测
决定保留或回退。最终代码的锁序、地址空间生命周期、页缓存代际、ext4 事务和双架构页表
逻辑均可由队员解释和现场修改，AI 不替代团队承担设计、验证和提交责任。

AI 生成内容曾出现错误或不完整判断。例如，单纯加强 VM 锁内存序只移动了卡死位置；扩大
VirtIO bounce、allocator 批处理和可写 ELF 共享在微基准中反而变慢；早期也曾把宿主资源
终止与 guest 故障混在同一轮长测里。团队通过 QEMU monitor、锁计数、紧邻 A/B、新鲜磁盘
覆盖层和官网日志发现这些问题，没有通过删除断言、伪造时间或修改评分脚本掩盖失败，而是
回退无收益候选并继续定位到无副作用等待谓词、COW 同步、目录 EOF 和产物启动 ABI 等根因。

本文同样由 AI 协助根据完整对话记录、\code{documents} 下的开发日志和 git 历史整理，
但所有数字均来自已有测试输出或官网结果，不由模型估算。队员复核了最终分支、提交边界、
日志标记和评分公式；官网成功结果是独立于 AI 的最终验证来源。

\section{可复现步骤}
\label{sec:reproduce}

\subsection{版本与源码检查}

从最终分支开始，并确认工作树没有混入诊断内核或旧磁盘产物：

\begin{Verbatim}[fontsize=\small,frame=single]
git fetch origin --prune
git switch finalsubmission
git rev-parse HEAD
git status --short
\end{Verbatim}

官网成功样本包含最后关键修复 \code{266bd1aa}。当前最终分支在此基础上增加
\code{8632a3e1} 的双架构确定性构建改进。若需要严格复核官网计时，应使用官网下载的
两份输出文件，而不是把后续本地构建时间与其混合。

\subsection{构建提交产物}

在项目要求的 Linux/Docker 工具链环境中执行：

\begin{Verbatim}[fontsize=\small,frame=single]
make all
file kernel-rv kernel-la
ls -lh kernel-rv kernel-la disk.img disk-la.img
\end{Verbatim}

\code{make all} 依次生成 RISC-V64 和 LoongArch64 内核及各自根盘。也可以分别检查：

\begin{Verbatim}[fontsize=\small,frame=single]
make kernel ARCH=riscv64
make kernel ARCH=loongarch64
cargo check -p os --target riscv64gc-unknown-none-elf
cargo check -p os --target loongarch64-unknown-none
\end{Verbatim}

LoongArch64 需要对应 C 交叉编译器来重建 lwext4；最终构建脚本会按 target 隔离静态库，
不能用 RISC-V64 的缓存产物替代。

\subsection{官方双盘运行}

把官方 BuildStorm 镜像作为 \code{TEST_DISK}，根盘使用当前架构生成的镜像。RISC-V64
最小复现方式为：

\begin{Verbatim}[fontsize=\small,frame=single]
make run-dual ARCH=riscv64 QEMU_SMP=8 QEMU_MEM=8G \
  TEST_DISK=/path/to/buildstorm-riscv64.img ROOT_DISK=disk.img
\end{Verbatim}

LoongArch64 使用：

\begin{Verbatim}[fontsize=\small,frame=single]
make run-dual ARCH=loongarch64 QEMU_SMP=8 QEMU_MEM=8G \
  TEST_DISK=/path/to/buildstorm-loongarch64.img ROOT_DISK=disk-la.img
\end{Verbatim}

若评测平台传入不同核数，以平台 QEMU 参数为准，judge 会读取最终 \code{cores} 字段。
最终验收不能只看 cargo 最后一行，必须同时出现以下四类证据：

\begin{Verbatim}[fontsize=\small,frame=single]
TOOLCHAIN_RESULT status=OK
MINIBUILD_RESULT status=OK
BUILDSTORM_RESULT ... status=OK ... bytes=<at least 512000> run=OK
#### OS COMP TEST GROUP END buildstorm-glibc ####
\end{Verbatim}

\subsection{短周期回归}

性能候选应先通过文件创建、并发 rename、truncate、稀疏文件、共享 inode、COW、
\code{MADV_DONTNEED}、线程组 wait、futex requeue、epoll ET、同核公平性和 8 核 SMP 压力，
再运行工具链和完整构建。诊断特性只用于统计，不应进入最终计时内核。每次 A/B 使用新鲜
覆盖层，防止 cargo 增量缓存把语义错误误判为加速。

\section{总结与后续方向}

ScintillaOS 最终解决的不是一个单点 syscall，而是一条跨越动态链接、线程组、futex、
用户内存、ext4、VirtIO、allocator、调度和 LoongArch MMU 的复杂构建链。性能上的最大收益
来自文件系统大粒度路径重构：哈希 CLOCK 块缓存、合并写回、固定顺序预读、inode 共享脏页
和 VFS 代际缓存；稳定性上的关键则是 COW 用户复制、线程组生命周期、无副作用 VM 锁等待
谓词、目录 EOF 和最终产物 ABI。

官网双架构工具链、minibuild、完整构建和产物启动全部通过，证明这些修改形成了可运行的
通用 Linux 兼容能力，而不是只让评分脚本输出成功。后续仍可在不改变工作量和计时语义的
前提下继续降低 RISC-V64 的 ext4 全卷串行区、改进块层异步队列和实现更成熟的页回写调度；
但任何新优化都应继续遵守本文的原则：先证明 Linux 语义，再做同条件 A/B，出现退化立即
回退，最后以双架构完整产物运行作为发布门槛。

\end{document}
